\SourcePage{307}{312}
\section{Invariant amplitudes}

In the helicity amplitudes a definite reference frame is used---the system of the centre of inertia. Meanwhile, in computing scattering amplitudes with the aid of the invariant perturbation theory (and also for the study of their general analytic properties) it is convenient to write the amplitudes in an explicitly invariant form.

If the particles participating in the reaction have no spin, then the amplitude of scattering depends only on invariant products of the 4-momenta of the particles. For a reaction of the type
\begin{equation}
a + b \to c + d
\tag{70.1}
\end{equation}
as these invariants one can choose any two of the three quantities $s$, $t$, $u$ defined in~\S\,66. Then the amplitude of scattering reduces to a single function $M_{fi} = f(s, t)$.

If the particles possess spins, then besides the kinematic invariants $s$, $t$, $u$ there also exist invariants which can be composed of the wave-function amplitudes of the particles (bispinors, 4-tensors, and so on). The amplitudes of scattering must then have the form
\begin{equation}
M_{fi} = \sum_n f_n(s, t)\,F_n,
\tag{70.2}
\end{equation}
where $F_n$ are invariants, linearly dependent on the wave-function amplitudes of all participating particles (and also on their 4-momenta). The coefficients $f_n(s, t)$ are called \textit{invariant amplitudes}.

Having chosen the wave amplitudes so that they correspond to particles with definite helicities, we obtain definite values of the invariants $F_n$. Then the helicity amplitudes of scattering are represented in the form of linear homogeneous combinations of invariant amplitudes $f_n$. From this it is clear that the number of independent functions $f_n$ coincides with the number of independent helicity amplitudes. Since the number of the latter is easily determined (as was explained in~\S\,69), thereby the task of constructing the invariants $F_n$ is also facilitated---we know in advance how many of them there should be.

Let us consider several examples. In all examples we shall assume that the interaction is $T$- and $P$-invariant; the latter property implies that the invariants $F_n$ must be true (and not pseudo) scalars.

\textbf{Scattering of a particle with spin 0 on a particle with spin 1/2.} For counting the number of invariants---or, what is the same, the number of elements of the matrix $S^J$ (i.e.\ the number of different sets of numbers $\lambda_1$, $\lambda_2$, $\lambda'_1$, $\lambda'_2$) --- in this case is equal to~4 ($\lambda_1 = \lambda'_1 = 0$, $\lambda_2$, $\lambda'_2 = \pm 1/2$). Taking into account $P$-invariance the number of independent elements reduces to two, after which the account of $T$-invariance no longer changes this number.

As the two independent invariants one can choose
\begin{equation}
F_1 = \bar{u}'u,\qquad F_2 = \bar{u}'\,(\gamma K)\,u.
\tag{70.3}
\end{equation}
Here $u = u(p)$, $u' = u(p')$ are the bispinor amplitudes of the initial and final fermions; $K = k + k'$, where $k$ and $k'$ are the 4-momenta of the initial and final bosons\,\footnote{At first glance one could also construct an invariant of the form $\bar{u}\,\sigma_{\mu\nu}\,k^\mu k'^\nu u$ (the matrices $\sigma_{\mu\nu}$ are defined in~(28.2)). It is easy, however, to verify its reducibility to the invariants~(70.3), if one takes into account the law of conservation $k' = p + k - p'$ and the equations
\[
(\gamma p)u = mu,\qquad \bar{u}'(\gamma p') = m\bar{u}',
\]
which the bispinor amplitudes satisfy.}.

$T$-invariance is obvious, noting that the products $\bar{u}'u$ and $\bar{u}'\gamma^\mu u$ are transformed under time reversal by the same law~(28.6), as the operators $\hat{\bar{\psi}}\hat{\psi}$ and $\hat{\bar{\psi}}\gamma^\mu\hat{\psi}$, the matrix elements of which they are: the product $\bar{u}'u$ is invariant by itself, and the 4-vector $\bar{u}'\gamma^\mu u$ is transformed by the law
\[
\bar{u}'\gamma^0 u \to \bar{u}'\gamma^0 u,\qquad \bar{u}'\boldsymbol{\gamma}u \to -\bar{u}'\boldsymbol{\gamma}u.
\]

\SourcePage{308}{313}
In the same manner the 4-momenta are transformed: $(K^0, \mathbf{K}) \to (K^0, -\mathbf{K})$, and the scalar product $F_2 = K_\mu(\bar{u}'\gamma^\mu u)$ is, consequently, invariant.

\textbf{Elastic scattering of two identical particles with spin 1/2.} For the counting of the number of independent helicity amplitudes it is convenient to start from the linear combinations of helicity states:
\[
\psi_{1g} = \psi_{++} + \psi_{--},\quad \psi_{2g} = \psi_{+-} - \psi_{-+},
\]
\[
\psi_{3g} = \psi_{+-} + \psi_{-+},\quad \psi_u = \psi_{+-} - \psi_{-+},
\]
where the indices ``$+$'', ``$-$'' indicate the values of the helicities ($\pm 1/2$) of the two particles. The states $1g$, $2g$, $3g$ are even, and the state $u$ is odd with respect to permutation of the particles. Therefore transitions $g \leftrightarrow u$ are forbidden, so that out of the original $16 - 6 = 10$ matrix elements there remain

\SourcePage{309}{314}
$P$ the functions $\psi_{1g}$, $\psi_{3g}$ and $\psi_{2g}$ have opposite parities; the prohibition of transitions between them reduces the number of independent amplitudes to six. Finally, $T$-invariance leads to the coincidence of amplitudes of transitions $1g \to 3g$ and $3g \to 1g$, so that there remain altogether five independent amplitudes. As the five independent invariants one can choose
\begin{equation}
\begin{aligned}
F_1 &= (\bar{u}'_1 u_1)(\bar{u}'_2 u_2),\qquad & F_2 &= (\bar{u}'_1\gamma^5 u_1)(\bar{u}'_2\gamma^5 u_2),\\[3pt]
F_3 &= (\bar{u}'_1\gamma^\mu u_1)(\bar{u}'_2\gamma_\mu u_2),\qquad & F_4 &= (\bar{u}'_1\gamma^\mu\gamma^5 u_1)(\bar{u}'_2\gamma_\mu\gamma^5 u_2)\\[3pt]
F_5 &= (\bar{u}'_1\sigma^{\mu\nu}u_1)(\bar{u}'_2\sigma_{\mu\nu}u_2),
\end{aligned}
\tag{70.4}
\end{equation}
where $u_1$, $u_2$ are the bispinor amplitudes of the initial, and $u'_1$, $u'_2$ those of the final particles. Permutation of the initial (or final) particles does not lead to new invariants: the new invariants are expressed through the old ones (see~\S\,28, Problem). But the expression~(70.2) with $F_n$ from~(70.4) does not explicitly take into account the requirement, according to which the permutation of two identical fermions must change the sign of the amplitude of scattering. An expression satisfying this requirement can be written in the form
\[
M_{fi} = [(\bar{u}'_1 u_1)(\bar{u}'_2 u_2)\,f_1(t, u) - (\bar{u}'_2 u_1)(\bar{u}'_1 u_2)\,f_1(u, t)] + \ldots
\tag{70.5}
\]
Under permutation $p'_1$ and $p'_2$ (or $p_1$ and $p_2$) the kinematic invariants $s \to s$, $u \to u$, $u \to t$, so that the indicated requirement is satisfied automatically.

\textbf{Elastic scattering of a photon on particles with spin 0 and 1/2.} The amplitude of these processes should conveniently be expressed with the help of unit spacelike 4-vectors $e^{(1)}$, $e^{(2)}$, satisfying the conditions
\begin{equation}
e^{(1)\,2} = e^{(2)\,2} = -1,\qquad e^{(1)}e^{(2)} = 0,
\tag{70.6}
\end{equation}
\[
e^{(1)}k = e^{(2)}k = 0,\qquad e^{(1)}k' = e^{(2)}k' = 0
\]
(for each of the two photons these 4-vectors can serve as the 4-orthonomals, with the help of which the invariant description of its polarisation properties is effected---see~\S\,8).

Let $k$ and $k'$ be the initial and the final 4-momenta of the photon, and $p$ and $p'$ the same for the scattering particle. Let us consider the 4-vectors
\begin{equation}
P^\lambda = p^\lambda + p'^\lambda - K^\lambda\frac{pK}{K^2},\qquad N^\lambda = e^{\lambda\mu\nu\rho}P_\mu q_\nu K_\rho,
\tag{70.7}
\end{equation}
where
\[
K = k + k',\qquad q = p - p' = k' - k.
\]
They are obviously mutually orthogonal. They are also orthogonal to the 4-vectors $K$, $q$, and hence to $k$, $k'$. Being orthogonal to the timelike 4-vector $K$ ($K^2 = 2kk' > 0$), they are spacelike (indeed, in the system of reference in which $\mathbf{K} = 0$, from $KP = 0$ it follows that $P_0 = 0$, and hence $P^2 < 0$). Having normalised $P$ and $N$, i.e.\ formed
\begin{equation}
e^{(1)\lambda} = \frac{N^\lambda}{\sqrt{-N^2}},\qquad e^{(2)\lambda} = \frac{P^\lambda}{\sqrt{-P^2}},
\tag{70.8}
\end{equation}
we obtain a pair of 4-vectors, possessing all the required properties. Let us note that $e^{(2)}$ is a true vector, and $e^{(1)}$ a pseudovector.

\SourcePage{310}{315}
Let us represent the amplitude of scattering of the photon in the form
\begin{equation}
M_{fi} = F^{\lambda\mu}\,e'^*_\lambda\,e_\mu,
\tag{70.9}
\end{equation}
having separated out the 4-vectors $e$ and $e'$ of the initial and the final photons.

The helicity of the photon takes only two values ($\pm 1$). Therefore for scattering of the photon on a particle with spin~0 the number of independent helicity amplitudes is the same as for the interaction of the particle with spin~0 and with spin~1/2, i.e.\ is equal to two. The tensor $F^{\lambda\mu}$ in~(70.9) must be constructed only from the 4-momenta of the particles. It can be represented in the form
\begin{equation}
F^{\lambda\mu} = f_1 e^{(1)\lambda}e^{(1)\mu} + f_2 e^{(2)\lambda}e^{(2)\mu},
\tag{70.10}
\end{equation}
where $f_1$, $f_2$ are the invariant amplitudes. Let us call attention to the fact that in $F^{\lambda\mu}$ there cannot be a term containing the product $e^{(1)\lambda}e^{(2)\mu}$, since this product is a pseudotensor and upon substitution into~(70.9) would give a pseudoscalar.

Finally, let us consider scattering of a photon on a particle with spin~1/2. For the counting of the number of independent helicity amplitudes we note that the total number of elements of the matrix $S^J$ in this case is~16 (the helicity of each of the two initial and two final particles takes two values). The requirement of $P$-invariance reduces this number to~8, after which the account of $T$-invariance brings it down to~6.

Let us represent the tensor $F_{\lambda\mu}$ in this case in the form
\begin{equation}
\begin{split}
F_{\lambda\mu} = G_0(e^{(1)}_\lambda e^{(1)}_\mu + e^{(2)}_\lambda e^{(2)}_\mu) + G_1(e^{(1)}_\lambda e^{(1)}_\mu - e^{(2)}_\lambda e^{(2)}_\mu) +{}\\[4pt]
{}+ G_2(e^{(1)}_\lambda e^{(2)}_\mu + e^{(2)}_\lambda e^{(1)}_\mu) + G_3(e^{(1)}_\lambda e^{(2)}_\mu - e^{(2)}_\lambda e^{(1)}_\mu),
\end{split}
\tag{70.11}
\end{equation}
where $G_0$, $G_3$ are true, and $G_1$, $G_2$ are pseudo-scalars. These and other bilinear quantities relative to the bispinor amplitudes of the fermions $\bar{u}(p')$ and $u(p)$, i.e.\ they have the form
\begin{equation}
G_n = \bar{u}(p')\,Q_n\,u(p).
\tag{70.12}
\end{equation}
The general form of the matrices (with respect to the bispinor indices) $Q_n$:
\begin{equation}
\begin{aligned}
Q_0 &= f_1 + f_2(\gamma K),\qquad & Q_1 &= \gamma^5[f_3 + f_4(\gamma K)],\\[3pt]
Q_2 &= \gamma^5[f_5 + f_6(\gamma K)],\qquad & Q_3 &= f_7 + f_8(\gamma K),
\end{aligned}
\tag{70.13}
\end{equation}
where $K = k + k'$. The coefficients $f_1, \ldots, f_8$ are invariant amplitudes, whose number has turned out to be~8 (instead of the required~6) because $T$-invariance has not yet been taken into account.

\SourcePage{311}{316}
Time reversal interchanges the initial and the final 4-momenta of the particles, also changing the signs of their spatial components:
\begin{equation}
(k_0, \mathbf{k}) \leftrightarrow (k'_0, -\mathbf{k}'),\qquad (p_0, \mathbf{p}) \leftrightarrow (p'_0, -\mathbf{p}');
\tag{70.14}
\end{equation}
the 4-vectors of polarisation of the photons are transformed according to
\begin{equation}
(e_0, \mathbf{e}) \leftrightarrow (e'^*_0, -\mathbf{e}'^*)
\tag{70.15}
\end{equation}
(cf.~(8.11a)), so that
\[
(e'^*_0 e_0,\; e'^*_i e_0,\; e'^*_0 e_k,\; e'^*_i e_k) \to (e'^*_0 e_0,\; -e'^*_0 e_i,\; e'^*_k e_i).
\]
By virtue of the last transformation the condition of invariance of the amplitude of scattering~(70.9) is equivalent to the requirement
\[
(F_{00}, F_{i0}, F_{ik}) \to (F_{00}, -F_{0i}, F_{ik}).
\]
On the other hand, as a consequence of the replacements~(70.14), we have
\[
(K_0, \mathbf{K}) \to (K_0, -\mathbf{K}),\quad (q_0, \mathbf{q}) \to (-q_0, \mathbf{q}),\quad (N_0, \mathbf{N}) \to (N_0, -\mathbf{N}),
\]
so that
\begin{equation}
(e^{(1,2)}_0, \mathbf{e}^{(1,2)}) \to (e^{(1,2)}_0, -\mathbf{e}^{(1,2)}).
\tag{70.16}
\end{equation}
From the expression~(70.11) it follows therefore that one should have
\[
G_{0,1,3} \to G_{0,1,3},\qquad G_2 \to -G_2.
\]
But under time reversal
\[
\bar{u}'\gamma^5 u \to -\bar{u}'\gamma^5 u,\qquad \bar{u}'\gamma^5(\gamma K)u \to \bar{u}'\gamma^5(\gamma K)u,
\]
as is clear from the laws of transformation of the pseudoscalar and pseudovector bilinear forms in~(28.6). Therefore from the expressions (70.12), (70.13) it is clear that by virtue of $T$-invariance of the amplitude of scattering it must be
\begin{equation}
f_3 = f_6 = 0.
\tag{70.17}
\end{equation}
